\documentclass[11pt]{letter}
\usepackage[margin=1in]{geometry}
\usepackage{times}
\usepackage[hidelinks]{hyperref}
\signature{Phuc Hao Do\\Da Nang Architecture University}
\address{Phuc Hao Do\\Da Nang Architecture University\\Da Nang, Vietnam\\haodp@dau.edu.vn}
\begin{document}
\begin{letter}{Editor-in-Chief\\IEEE Transactions on Cognitive Communications and Networking}
\opening{Dear Editor,}

I am pleased to submit the manuscript ``When Do Kolmogorov--Arnold Networks Help in Wireless?
A Matched-Capacity, Non-Saturated Benchmark'' for consideration as a research paper in the IEEE
Transactions on Cognitive Communications and Networking.

Kolmogorov--Arnold networks (KANs) are being adopted rapidly across the wireless physical layer,
usually with two claims: that they match or beat multilayer perceptrons at fewer parameters, and
that their learned edge functions make them interpretable. These claims are typically demonstrated
in a showcase style, against an unmatched baseline and on a possibly saturated benchmark. Our paper
asks, and answers by measurement, a more useful question: when does a KAN actually help in wireless?

The contribution is a controlled, reproducible benchmark rather than a new architecture. We
introduce a parameter-matched, multi-seed protocol and apply it across three non-saturated regimes,
power-amplifier behavioral modeling, OFDM channel estimation, and modulation classification on the
real RML2016.10a dataset, each with a strong classical or structural baseline (generalized memory
polynomial, LMMSE, and a convolutional network). The result is deliberately two-sided and honest:
KANs are competitive and win in specific regimes (high-SNR classification, low-data, richer bases),
but they never dominate, are beaten by the right classical or structural method, carry an
interpretability cost of more than a decibel that we measure directly, and run four to thirty times
slower than a matched MLP. We distill the findings into a regime map and a six-point reporting
checklist for the community, and we release all code, data-preparation scripts, and a
single-source-of-truth claims file at \url{https://github.com/haodpsut/kan-wireless-benchmark}, so
that every number in the paper can be regenerated and verified.

We believe the work fits the scope of the Transactions squarely: it concerns machine learning for
cognitive and adaptive wireless systems, and it provides a methodology that helps the community
evaluate such models honestly. To our knowledge the wireless-specific question has not been studied
under matched capacity before; the closest prior work is a general-machine-learning fairness study,
which we cite and extend to the wireless setting.

We confirm that this manuscript is original, has not been published elsewhere, and is not under
review at any other venue. A separate paper of ours applies KANs to a different problem (adaptive
control in satellite quantum key distribution) and is under review elsewhere; it shares no results,
datasets, or claims with this benchmark, which is an independent methodological contribution. The
author has no conflicts of interest to declare.

Thank you for your consideration.

\closing{Sincerely,}
\end{letter}
\end{document}
